\section{Conclusion and Future Work}\label{sec:conclusion}
This paper presented a high-throughput, GPU-accelerated C++/CUDA framework for real-time 5D image and video segmentation on a joint color-spatial manifold. By optimizing grid occupancy, utilizing on-chip shared memory atomics, and staging memory transfers through pinned DMA channels, the pipeline achieves theoretical throughputs up to $750\text{ FPS}$ on an RTX 5070 Ti, operating with significant headroom under the $120\text{ FPS}$ ($8.33\text{ ms}$) budget constraint.

The quantum-inspired engine emulates the multi-qubit state overlap of the quantum SWAP test by mapping pixel coordinates to phase states. Due to the monotonic nature of the cosine-squared mapping over the $[0, \pi/4]$ phase range, the quantum-inspired backend achieves visual and mathematical equivalence to the classical Euclidean solver under varying illumination conditions. Leveraging SFU hardware-accelerated math intrinsics (\texttt{\_\_cosf()}) bypasses standard ALU emulation, maintaining execution latency parity with classical Euclidean K-Means.

Future work will focus on integrating spatiotemporal constraints (such as Hungarian association or optical flow) to stabilize boundary jitter in continuous video streams. Additionally, we plan to port the phase-space similarity metrics to native NISQ hardware to evaluate real state-preparation and measurement delays under hardware noise, and optimize execution profiles on low-wattage embedded architectures.
